
Semantic entropy outperforms token entropy by 0.155 AUROC on TriviaQA --- then loses by 0.066 AUROC on TruthfulQA. Same model, opposite orderings. This reversal is not noise: it exposes a task-structure condition that determines when semantic-level uncertainty estimation helps and when it fails. We evaluate four major uncertainty proxies --- semantic entropy (SE), token entropy (TE), SelfCheckGPT BERTScore (SCG), and verbalized confidence (VC) --- at Llama-2-7B scale under identical conditions and find that SE outperforms TE by +0.155 AUROC on TriviaQA (N=98; bootstrap 95\% CI on the gap excludes zero; per-method CIs marginally overlap by 0.063), with the gap mechanistically explained by TE aggregating surface-form variation within semantically equivalent paraphrase clusters (mean intra-cluster TE variance = 7.152 $nats^2, 71\times$ the 0.1 $nats^2$ significance threshold). This advantage reverses on TruthfulQA (N=141; TE = 0.511 > SE = 0.445), where adversarial misconceptions produce deterministic-wrong outputs that SE's paraphrase-filtering mechanism cannot handle. We additionally show that BERTScore-based SCG is not equivalent to SE on short factual QA (raw |SCG $-$ SE| = 0.092; corrected delta using consistent AUROC orientation = 0.336; gate threshold = 0.03), and that verbalized confidence at 7B scale produces a near-degenerate confidence distribution (ECE = 0.430, 5 distinct values). These findings establish the task-structure condition for SE's superiority: SE should be preferred when incorrect model outputs are paraphrase-diverse, and TE when they are deterministically wrong --- a scope condition not identified in prior 65B-scale evaluations.

